% Topic 1.2.08 — Mean-Square Ergodic Processes.

\section{Mean-Square Ergodic Processes}
\label{sec:1.2.08-ergodicity}

\begin{ticket}
Mean-square ergodic stochastic processes (in expectation, variance, correlation function). Criterion and sufficient condition for mean-square ergodicity in expectation.
\end{ticket}

\subsection*{Theory}

The strong law of large numbers (\cref{thm:slln}) says that the
time average of an i.i.d.\ sequence converges to the ensemble mean.
For dependent processes, the analogous question is: when does the
time average over a single trajectory recover the population
parameters? The mean-square framework gives the cleanest answer
under second-moment hypotheses.

\medskip
\noindent\textbf{Wide-sense stationary processes.}

\begin{definition}[WSS process]
\label{def:wss}
A stochastic process $X = (X_t)_{t \in T}$ with $T \in \{\Z, \R, [0, \infty)\}$
and $X_t \in L^2$ for all~$t$ is \emph{wide-sense stationary
(WSS)} if
\begin{enumerate}[label=(\alph*)]
  \item $\E X_t = m$ is constant in~$t$;
  \item $\Cov(X_s, X_t)$ depends only on the time lag $\tau = t - s$.
\end{enumerate}
The function $K(\tau) = \Cov(X_t, X_{t + \tau})$ is the
\emph{(auto-)covariance function} (or \emph{correlation function} in
parts of the Russian literature, when the process is centred).
\end{definition}

By \cref{prop:cov-properties}, $K(0) = \Var X_t \geq 0$ and $K(-\tau) = K(\tau)$;
moreover $K$ is positive-semidefinite as a kernel, so it admits a
spectral representation \cite[Ch.~VI, \S~1]{Sevastyanov1982}.

\medskip
\noindent\textbf{Mean-square ergodicity.} For a WSS process, the
time-average estimator of the mean is
\[
  \bar X_T \;=\; \frac{1}{T} \int_0^{T} X_t\, dt \quad \text{(continuous time)}, \qquad
  \bar X_n \;=\; \frac{1}{n} \sum_{k=1}^{n} X_k \quad \text{(discrete time)}.
\]
We treat the continuous-time case below; the discrete case is
parallel with sums in place of integrals.

\begin{definition}[Mean-square ergodicity]
\label{def:ms-ergodic}
A WSS process~$X$ is \emph{mean-square ergodic in the mean} if
$\bar X_T \to m$ in mean square as $T \to \infty$, i.e.\
$\E (\bar X_T - m)^2 \to 0$. Analogously, $X$ is mean-square
ergodic \emph{in variance} (resp.\ \emph{in correlation function})
if the time-average estimators of $\Var X_t$ (resp.\ of
$K(\tau)$ for fixed lag~$\tau$) converge in $L^2$ to the
corresponding population quantities.
\end{definition}

\begin{theorem}[Criterion for mean-square ergodicity in the mean]
\label{thm:ms-ergodic-criterion}
Let $X$ be a WSS process with mean~$m$ and covariance function~$K$.
Then $X$ is mean-square ergodic in the mean if and only if
\[
  \frac{1}{T^2} \int_0^{T} \int_0^{T} K(t - s)\, ds\, dt \;\to\; 0 \qquad (T \to \infty).
\]
Equivalently, after a change of variables $\tau = t - s$, this is
\[
  \frac{2}{T} \int_0^{T} \Bigl(1 - \frac{\tau}{T}\Bigr) K(\tau)\, d\tau \;\to\; 0.
\]
\end{theorem}
\proofof{ms-ergodic-criterion}

\begin{corollary}[Sufficient condition]
\label{cor:ms-ergodic-sufficient}
If $K(\tau) \to 0$ as $\tau \to \infty$, then $X$ is mean-square
ergodic in the mean. In particular, this holds whenever
$\int_0^{\infty} |K(\tau)|\, d\tau < \infty$.
\end{corollary}
\proofof{ms-ergodic-sufficient}

\begin{remark}
The criterion isolates the role of long-range correlation. If
$K(\tau)$ decays to zero, the variance of the time average vanishes
at rate $O(1/T)$ (or better) — sample paths self-average. If
$K(\tau)$ persists at a positive level — for instance, if
$X_t = Y$ for a single random variable~$Y$, in which case
$K(\tau) \equiv \Var Y$ — the time average concentrates not around
$m$ but around the random variable~$Y$, and ergodicity fails.
Mean-square ergodicity in variance and in correlation is reduced to
the same criterion applied to the auxiliary processes
$(X_t - m)^2$ and $(X_t - m)(X_{t + \tau} - m)$, provided the
necessary fourth moments exist
\cite[Ch.~VI, \S~5]{Sevastyanov1982}.
\end{remark}

\subsection*{Proofs}

\begin{proof}[\Cref{thm:ms-ergodic-criterion}]
\proofanchor{ms-ergodic-criterion}
By Fubini and linearity of expectation
(\cref{prop:exp-properties}\,(a)),
$\E \bar X_T = \frac{1}{T} \int_0^T \E X_t\, dt = m$. Hence
\[
  \E (\bar X_T - m)^2 = \Var \bar X_T = \frac{1}{T^2} \int_0^T \!\!\int_0^T \Cov(X_s, X_t)\, ds\, dt = \frac{1}{T^2} \int_0^T \!\!\int_0^T K(t - s)\, ds\, dt,
\]
using stationarity in the last step. So $L^2$-convergence
$\bar X_T \to m$ is equivalent to the displayed double integral
tending to~$0$, which is condition~(\ref{thm:ms-ergodic-criterion}).

For the equivalent single-integral form, change variables to
$\tau = t - s$, $\sigma = s$ on the square
$[0, T] \times [0, T]$. The image is
$\{(\tau, \sigma) : 0 \leq \sigma \leq T,\ -\sigma \leq \tau \leq T - \sigma\}$;
fix $\tau$ and integrate over $\sigma \in [\max(0, -\tau), \min(T, T - \tau)]$,
which has length $T - |\tau|$ for $|\tau| \leq T$. By symmetry
$K(-\tau) = K(\tau)$, the integral becomes
\[
  \frac{1}{T^2} \int_{-T}^{T} (T - |\tau|) K(\tau)\, d\tau = \frac{2}{T} \int_0^T \Bigl(1 - \frac{\tau}{T}\Bigr) K(\tau)\, d\tau. \qedhere
\]
\end{proof}

\begin{proof}[\Cref{cor:ms-ergodic-sufficient}]
\proofanchor{ms-ergodic-sufficient}
Suppose $K(\tau) \to 0$ as $\tau \to \infty$, so given $\varepsilon > 0$
there is~$M$ with $|K(\tau)| < \varepsilon$ for $\tau \geq M$. For
$T > M$, split the inner integral in
\cref{thm:ms-ergodic-criterion}:
\[
  \frac{2}{T} \int_0^T \Bigl(1 - \frac{\tau}{T}\Bigr) K(\tau)\, d\tau = \frac{2}{T} \int_0^M (\cdot)\, d\tau + \frac{2}{T} \int_M^T (\cdot)\, d\tau.
\]
The first integral is bounded by $2M \cdot \max_{[0, M]} |K|/T \to 0$
as $T \to \infty$. The second is bounded by
$\frac{2}{T} \int_M^T (1 - \tau/T) \varepsilon\, d\tau \leq \varepsilon$
(since the weight $1 - \tau/T \leq 1$ and the interval has
length~$\leq T$). Hence the lim sup is $\leq \varepsilon$ for every
$\varepsilon > 0$, i.e.\ the integral tends to~$0$, and
\cref{thm:ms-ergodic-criterion} gives mean-square ergodicity. The
absolute-integrability case $\int_0^{\infty} |K| < \infty$ is even
simpler: $\frac{2}{T} \int_0^T (1 - \tau/T) |K(\tau)|\, d\tau \leq \frac{2}{T} \int_0^{\infty} |K|\, d\tau \to 0$.
\end{proof}

\subsection*{Examples}

\begin{example}[White noise]
\label{ex:white-noise}
Let $(X_n)_{n \in \Z}$ be i.i.d.\ with mean~$m$ and variance
$\sigma^2$. Then $K(0) = \sigma^2$ and $K(\tau) = 0$ for
$\tau \neq 0$. Trivially $K(\tau) \to 0$, so by
\cref{cor:ms-ergodic-sufficient} the process is mean-square
ergodic in the mean. The discrete-time analogue of the variance
of the average is $\Var \bar X_n = \sigma^2/n \to 0$ — exactly
\cref{ex:binomial-moments} re-cast.
\end{example}

\begin{example}[Ornstein–Uhlenbeck process]
\label{ex:ou-ergodic}
The continuous-time Gaussian process from \cref{prob:fdd-marginal}
with $K(\tau) = e^{-|\tau|}$ satisfies $K(\tau) \to 0$
exponentially, so it is mean-square ergodic in the mean. The
sufficient condition's quantitative form gives the rate
\[
  \Var \bar X_T = \frac{2}{T} \int_0^T \bigl(1 - \tau/T\bigr) e^{-\tau}\, d\tau = \frac{2}{T}\bigl(1 - (1 - e^{-T})/T\bigr) \sim \frac{2}{T},
\]
asymptotically the same $1/T$ rate as for white noise — long-range
correlations are weak enough that they do not slow self-averaging
beyond a constant factor.
\end{example}

\begin{example}[Constant random process — non-ergodic]
\label{ex:constant-process}
Let $Y$ be any random variable with $\Var Y = \sigma^2 > 0$ and
set $X_t \equiv Y$ for all~$t$. Then $\E X_t = \E Y = m$ and
$K(\tau) = \Cov(X_s, X_{s + \tau}) = \Var Y = \sigma^2$ for every
lag~$\tau$. The criterion fails because
$\frac{1}{T^2} \int_0^T \int_0^T K(t - s)\, ds\, dt = \sigma^2 \neq 0$.
Indeed, $\bar X_T = Y$ for every~$T$, so the time average is
identically $Y$ and concentrates at the random variable~$Y$
rather than at the constant~$m$.
\end{example}

\begin{problem}
\label{prob:periodic-ergodic}
Let $X_t = A \cos(\omega t + \Phi)$ where $A, \omega \in \R$ are
constants and $\Phi \sim \mathrm{Uniform}(0, 2\pi)$. Show that $X$
is WSS, compute its correlation function, and decide whether $X$
is mean-square ergodic in the mean.
\end{problem}

\begin{proof}[Solution]
$\E X_t = A\, \E \cos(\omega t + \Phi) = \frac{A}{2\pi} \int_0^{2\pi} \cos(\omega t + \varphi)\, d\varphi = 0$
since $\cos$ integrates to $0$ over a full period. For the
covariance,
\[
  \E (X_s X_t) = A^2 \E \cos(\omega s + \Phi)\cos(\omega t + \Phi) = \tfrac{A^2}{2}\, \E\bigl(\cos(\omega(t - s)) + \cos(\omega(s + t) + 2\Phi)\bigr).
\]
The second term vanishes after taking expectation (its phase
$2\Phi$ is uniform on $(0, 4\pi)$), leaving
$\Cov(X_s, X_t) = K(t - s) = \frac{A^2}{2} \cos(\omega(t - s))$.
Both conditions of \cref{def:wss} hold (mean is the constant~$0$
and covariance depends only on the lag), so~$X$ is WSS.

For ergodicity, $K(\tau) = \frac{A^2}{2} \cos(\omega \tau)$ does
\emph{not} tend to $0$, so the sufficient condition fails. Apply
\cref{thm:ms-ergodic-criterion} directly: for $\omega \neq 0$,
\[
  \frac{2}{T} \int_0^T \bigl(1 - \tau/T\bigr) \tfrac{A^2}{2} \cos(\omega \tau)\, d\tau = \frac{A^2}{T} \cdot \frac{1 - \cos(\omega T)}{\omega^2 T} = O(1/T^2) \to 0,
\]
so the process \emph{is} mean-square ergodic in the mean despite
the non-decaying~$K$ — the oscillations average to zero. (For
$\omega = 0$ we recover \cref{ex:constant-process} and ergodicity
fails.) The takeaway: the criterion of
\cref{thm:ms-ergodic-criterion} is more delicate than the
sufficient condition of \cref{cor:ms-ergodic-sufficient}.
\end{proof}
